% Standalone problem document, generated from the adjacent README.md.
% Compile with XeLaTeX. Mathematical content is maintained in Markdown.
\documentclass[11pt,a4paper]{article}
\usepackage[fontsize=10.5pt]{scrextend}
\usepackage[margin=22mm,headheight=15pt,headsep=9mm,footskip=11mm]{geometry}
\usepackage{fontspec}
\setmainfont{texgyrepagella}[Extension=.otf,UprightFont=*-regular,BoldFont=*-bold,ItalicFont=*-italic,BoldItalicFont=*-bolditalic]
\setsansfont{texgyreheros}[Extension=.otf,UprightFont=*-regular,BoldFont=*-bold,ItalicFont=*-italic,BoldItalicFont=*-bolditalic]
\setmonofont{texgyrecursor}[Extension=.otf,UprightFont=*-regular,BoldFont=*-bold,ItalicFont=*-italic,BoldItalicFont=*-bolditalic,Scale=MatchLowercase]
\usepackage{amsmath,amssymb}
\usepackage{unicode-math}
\setmathfont{texgyrepagella-math.otf}
\usepackage{xcolor,microtype,enumitem,needspace,fancyhdr}
\usepackage{bookmark}
\definecolor{ink}{HTML}{17344B}
\definecolor{accent}{HTML}{197D80}
\definecolor{muted}{HTML}{596773}
\hypersetup{colorlinks=true,linkcolor=accent,urlcolor=accent,pdftitle={IV-02 - Exact
determinant ranges of general tridiagonal interval
matrices},pdfauthor={Open Problems in Numerical Linear Algebra}}
\urlstyle{same}
\setlength{\parindent}{0pt}
\setlength{\parskip}{5pt plus 1pt minus 1pt}
\linespread{1.04}
\setlength{\emergencystretch}{3em}
\widowpenalty=10000
\clubpenalty=10000
\displaywidowpenalty=10000
\allowdisplaybreaks[1]
\setlist{nosep,leftmargin=1.5em}
\providecommand{\tightlist}{\setlength{\itemsep}{0pt}\setlength{\parskip}{0pt}}
\providecommand{\pandocbounded}[1]{#1}
\pagestyle{fancy}
\fancyhf{}
\fancyhead[L]{\sffamily\footnotesize\color{muted}OPEN PROBLEMS IN NUMERICAL LINEAR ALGEBRA}
\fancyhead[R]{\sffamily\bfseries\footnotesize\color{accent}IV-02}
\fancyfoot[L]{\parbox[b]{0.9\textwidth}{\sffamily\fontsize{8}{10}\selectfont\color{muted}Editorial ratings; a bounded search cannot certify that a problem remains unsolved.\\Literature check: 2026-09-08}}
\fancyfoot[R]{\sffamily\footnotesize\color{muted}\thepage}
\renewcommand{\headrulewidth}{0.3pt}
\renewcommand{\headrule}{\hbox to\headwidth{\color{accent}\leaders\hrule height \headrulewidth\hfill}}
\makeatletter
\renewcommand{\section}{\Needspace{5\baselineskip}\@startsection{section}{1}{0pt}{12pt plus 2pt minus 2pt}{4pt}{\sffamily\bfseries\large\color{ink}}}
\renewcommand{\subsection}{\Needspace{4\baselineskip}\@startsection{subsection}{2}{0pt}{9pt plus 1pt minus 1pt}{3pt}{\sffamily\bfseries\normalsize\color{ink}}}
\makeatother
\setcounter{secnumdepth}{0}
\begin{document}
{\sffamily\small\color{muted}Intervals and absolute value equations\par}
\vspace{3pt}
{\raggedright\hyphenpenalty=10000\exhyphenpenalty=10000\sffamily\bfseries\fontsize{21}{24}\selectfont\color{ink}Exact
determinant ranges of general tridiagonal interval matrices\par}
\vspace{7pt}
{\sffamily\small\textbf{Difficulty:} challenging\quad\textbf{Importance:} interesting
to the community\par}
\vspace{4pt}
{\color{accent}\hrule height 0.8pt}
\vspace{6pt}
\textbf{Topic:} interval linear algebra; determinants; computational
complexity

\textbf{Status:} open; admitted after a bounded literature search found
no resolution.

\subsection{Problem statement}\label{problem-statement}

For \(n\geq2\), let the input consist of \(3n-2\) closed real intervals
with rational endpoints: diagonal intervals
\([\underline a_i,\overline a_i]\) for \(1\leq i\leq n\), upper-diagonal
intervals \([\underline b_i,\overline b_i]\), and lower-diagonal
intervals \([\underline c_i,\overline c_i]\) for \(1\leq i<n\). All
lower endpoints are at most their upper endpoints. Define

\[
\mathcal T=\{T\in\mathbb R^{n\times n}:
T_{ii}\in[\underline a_i,\overline a_i],\quad
T_{i,i+1}\in[\underline b_i,\overline b_i],\quad
T_{i+1,i}\in[\underline c_i,\overline c_i],\quad
T_{ij}=0\text{ if }|i-j|>1\}.
\]

All uncertain entries vary independently. Does a deterministic algorithm
exist that returns the two exact rational numbers

\[
d_- = \min_{T\in\mathcal T}\det T,
\qquad
d_+ = \max_{T\in\mathcal T}\det T
\]

in time polynomial in the total binary input length? The target is the
exact range \([d_-,d_+]\), including instances containing singular
matrices and intervals crossing zero. Rational output is appropriate
because the determinant is affine in each entry separately and its
extrema are attained at endpoint matrices. The polynomial bound must be
uniform in \(n\).

\subsection{References}\label{references}

Horáček, Hladík, and Matějka,
\href{https://doi.org/10.13001/1081-3810.3719}{\emph{Determinants of
interval matrices}}, Electronic Journal of Linear Algebra 33 (2018),
99--112, \textbf{§5.4, p.~106}, immediately after Proposition 5.6;
\href{https://journals.uwyo.edu/index.php/ela/article/download/1831/1831/1831}{journal
PDF}. The corresponding location in
\href{https://arxiv.org/abs/1809.03736v1}{arXiv:1809.03736v1} is §6.4,
p.~9, after Proposition 6.7. The source proves polynomial computability
for interval tridiagonal H-matrices and explicitly leaves general
tridiagonal interval matrices open.

\subsection{Status check (2026-09-08)}\label{status-check-2026-09-08}

Searches for \textit{tridiagonal interval determinant complexity},
\textit{tridiagonal determinant range}, and
\textit{tridiagonal interval determinant polynomial 2026} found no
resolution for independent-entry exact ranges. Two superficially
conflicting Thirupathi--Thamaraiselvan papers,
\href{https://doi.org/10.28924/2291-8639-21-2023-20}{\emph{Symbolic
Algorithm for Inverting General k-Tridiagonal Interval Matrices}} and
\href{https://doi.org/10.28924/2291-8639-21-2023-87}{\emph{A Symbolic
Algorithm for Solving Doubly Bordered k-Tridiagonal Interval Linear
Systems}}, both published in 2023, use different generalized arithmetic
(§2.2, pp.~3--4). Their multiplication centers the result at the product
of midpoints and takes the smaller distance to the standard product
endpoints. Consequently it sends \([1,2]\) and \([2,3]\) to
\([2,11/2]\), while independent products range over \([2,6]\). Their
symbolic determinant algorithms therefore do not establish exact ranges
in the sense defined above.
\end{document}
